\section{Literature Review}
\label{sec:litreview}
This section synthesises twenty journal publications from 2021--2025: ten review articles that map the field and ten original studies that test specific detectors. They are grouped by theme rather than listed one by one, because the aim is to show what the evidence establishes about representations, evaluation and deployment, and where it stops.

\subsection{How the Field Frames the Problem}
The broad surveys agree that fake-review detection is several tasks rather than one: classifying individual review texts, identifying suspicious reviewers or groups, and ranking products or accounts for moderation \cite{ref13,ref8,ref12}. Mohawesh et al.\ document the move from handcrafted features to learned text representations and benchmark neural and transformer models, while naming reliable labels and comparable evaluation as open problems \cite{ref13}. He et al.\ separate signals that originate in the text from those derived from users, products and their links, which implies that every comparison should state what information is available at prediction time \cite{ref8}. Mewada and Dewang organise the literature around two operations, designing indicators of deception and designing a recognition method, and identify large, reliably labelled datasets as the main bottleneck \cite{ref12}. Paul and Nikolaev place opinion spam within organised, adversarial manipulation on e-commerce platforms and explain why platform and regulatory responses have had limited success \cite{ref16}. Ennaouri and Zellou focus on machine-learning approaches built on review content, sentiment and feature selection, and stress that reviews arrive continuously, which makes automation and processing efficiency part of the problem rather than an implementation detail \cite{ref5}. Taken together, these surveys define the target of detection (a review, a reviewer or a group), the evidence available at prediction time, and the label source as the three choices that any experiment must state before its results can be interpreted. A broad review cannot itself establish how well a specific new classifier performs, so these papers frame the design of the present study rather than supply benchmark values for it.

Three reviews widen the frame beyond classification accuracy. Walther et al.\ synthesise 15 consumer studies and group the cues people use to judge authenticity into review, text, reviewer, seller and platform characteristics; human judgement is therefore multi-cue, not purely lexical \cite{ref19}. Lim et al.\ analyse 229 studies with structural topic modelling and link detection to the decisions of firms, platforms and consumers \cite{ref9}. Gupta et al.\ attribute weak generalisation to limited training data, variation across datasets, concept drift and new fabrication methods \cite{ref7}. Duma et al.\ add the gaps of low-resource and code-mixed languages and of representations that ignore correlations between ratings, aspects and emotions \cite{ref4}, and their later review of graph neural networks (GNNs) covers work up to January 2024, showing that relational patterns can expose spam that looks normal in isolation while graph construction remains unsettled \cite{ref3}.

\subsection{Text-Based Detectors}
Contextual language models dominate recent text-only work. Lu et al.\ combine BERT representations, sentiment knowledge (SKEP) and a TextCNN in the BSTC model and report accuracy of 93.44\%, 91.25\% and 92.86\% on hotel, restaurant and doctor datasets \cite{ref11}. Mohawesh et al.\ couple RoBERTa with an LSTM, explain predictions with SHAP and attention, and report 96.03\% accuracy on OpSpam and 93.15\% on a second deception corpus \cite{ref15}. Geetha et al.\ tune DeBERTa with Monarch Butterfly Optimisation and test noise and augmentation; their accuracy ranges from 78\% on Amazon to 98\% on a fake-review dataset and 91\% on Deceptive Opinion Spam \cite{ref6}. Zhao et al.\ build a Chinese physician-review corpus on which BERT reaches 98.36\% precision and an F2 score of 97.97\%, ahead of traditional models, but they note that the corpus is small, single-platform and oversampled \cite{ref20}. Salminen et al.\ generate product reviews with GPT-2 and show that a RoBERTa detector reaches about 0.97 F1 on machine-generated text, but that this setting is not interchangeable with naturally deceptive human reviews \cite{ref17}.

These studies share strengths and weaknesses. Most test more than one dataset or domain, and Mohawesh et al.\ make interpretability an explicit goal through SHAP and attention \cite{ref15}. However, the wide accuracy range reported by Geetha et al.\ for a single architecture shows that a percentage obtained on one corpus says little about another \cite{ref6}, and Zhao et al.\ caution that oversampling and a single platform restrict inference \cite{ref20}. GPT-2-era generated reviews also do not represent later large language models or paraphrasing attacks \cite{ref17}. None of these text-only studies reports how its detector behaves when the same content is presented in a different surface form, and none compares the transformer with an unmodified linear baseline under an identical, leakage-controlled partition. Both omissions are directly relevant to deployment, where input formatting is not controlled and where the cheaper model is the default unless the more expensive one is shown to be better.

\subsection{Behavioural, Multi-View and Imbalance-Aware Detectors}
A second line of work adds evidence that a text classifier cannot see. Liu et al.\ jointly model text, users, products and fine-grained aspects and report AUC values of 86.65\%, 91.89\% and 93.26\% on YelpChi, YelpNYC and YelpZIP \cite{ref10}. The EUPHORIA multi-view network of Andresini et al.\ combines content with behavioural features and reaches AUC-ROC 0.813 on Yelp hotels but only 0.708 on restaurants \cite{ref1}, a gap that warns against reading any single headline score as general. Mohawesh et al.\ integrate review, reviewer and product information in an explainable ensemble and report AUC gains of one to seven percentage points over their baselines \cite{ref14}. Budhi et al.\ build 133 content and behaviour features and show that 99.3\% accuracy on genuine reviews can coexist with 7.73\% on the fake class; resampling and parallel cross-validation raise their best accuracy to about 89\% \cite{ref2}. Singhal and Kashef show that sampling and weighted stacking of conventional classifiers can be competitive, so gains do not come only from larger models \cite{ref18}.

The multi-view studies demonstrate that behavioural and relational evidence adds information, and their use of several real Yelp collections is a strength. Their requirements are also their limitation: user and product identifiers must exist, be reliable and transfer to new platforms, and the platform-provided labels they learn from are themselves the output of a proprietary filter. The imbalance studies make a complementary methodological point. When fake reviews are rare, a single accuracy figure can hide a classifier that almost never detects them \cite{ref2}, so class-specific precision and recall and a macro-averaged score are necessary. Budhi et al.\ also tie detection to big-data practice by parallelising cross-validation, but processing many records is not the same as demonstrating throughput, memory use or latency at platform scale.

\subsection{Datasets and Label Sources}
Across the twenty studies, four kinds of label appear, and they define different tasks. Platform-filter labels, as in the YelpChi, YelpNYC and YelpZIP collections and the Yelp hotel and restaurant data, record what a commercial filter hid from readers \cite{ref10,ref1}; they are large and realistic but are proxies whose errors are unknown. Crowdsourced labels, as in the Deceptive Opinion Spam corpus used by several transformer studies \cite{ref15,ref6}, are certain by construction but describe reviews written on request by people who never stayed at the hotel. Machine-generated labels, as in the GPT-2 reviews of Salminen et al., describe a third population again \cite{ref17}, and purpose-built domain corpora such as the physician reviews of Zhao et al.\ add a fourth, with its own collection rules \cite{ref20}. The surveys repeatedly identify reliable labels as the field's main bottleneck \cite{ref12,ref13,ref7}. A score is therefore comparable only with scores obtained on the same label type and partition, which is why the present study compares its models within one corpus and one set of partitions rather than against published figures from other corpora.

\subsection{Synthesis and Research Gaps}
The reviewed evidence supports four conclusions. First, contextual transformers are the strongest text-only models in these studies \cite{ref11,ref15,ref6}, but their reported scores come from different corpora, splits and metrics (accuracy, AUC, F1, F2), so they cannot be ranked against one another or against a simple baseline. Second, the reviewed original studies do not pair a transformer with a fixed linear baseline on one leakage-controlled partition and test the difference with a paired statistical procedure, which is the comparison needed to show that a transformer's gain is real and worth its cost. Third, multi-view and graph models gain from behavioural data that many corpora lack \cite{ref1,ref10,ref14,ref3}, so text-only performance remains the relevant reference when such data are missing. Fourth, generalisation across platforms, time and generators is repeatedly named as the central unsolved issue \cite{ref7,ref4,ref17}, yet most experiments still use random review-level splits. Finally, although Budhi et al.\ parallelise cross-validation \cite{ref2} and Geetha et al.\ examine training and inference costs \cite{ref6}, the studies do not provide a standardised report of throughput, memory and training time alongside predictive quality, even though continuous review streams are a stated motivation \cite{ref5}. These gaps define the design of the present study: an identical, leakage-controlled partition for a linear baseline and a transformer, paired uncertainty estimates, robustness checks and measured computational cost. The literature thus treats fake-review detection as a data and evaluation problem as much as a modelling problem: a credible detector must be assessed for what its labels represent, where it works, how stable its advantage is and what resources it needs.
